% Leading result derived before conditional2h outcomes; explicit-constant
% extension derived before reading observable1h outcomes. No new trajectories.
\section{First-Order Regime Value Versus Second-Order State Value}
The state-value ceiling does not rule out a substantial benefit from identifying
the covariance regime: that information can change the leading variance
coefficient of a resource program.
The distinction can be made exact against every affordable fixed training
prefix, rather than only against programs that spend the whole budget.

Consider the running mean with a fixed public initial-error radius $r>0$,
finite positive participation and spacing sets, bounded delay $D$, and
$0\le\lambda\le1-\gamma<1$.
Write the message and clock budgets as $\lfloor\beta_m B\rfloor$ and
$\lfloor\beta_e B\rfloor$, where $\beta_m,\beta_e>0$.
For program $a=(q,b)$ define
\begin{align*}
a_a&=\min\{\beta_m/(h+q),\beta_e/b\},\\
v_j(a)&=q^{-1}+\theta_j\frac{1+\lambda^b}{1-\lambda^b},\\
K_j(a)&=v_j(a)/a_a.
\end{align*}
All programs retain the same running-mean rule; a fixed alternative may choose
any deterministic scheduled prefix, including zero updates, within its budgets.
Let
\begin{align*}
K_\pi&=\min_a\sum_j\pi_jK_j(a),\\
a_j^*&\in\arg\min_aK_j(a),\\
\Delta&=K_\pi-\sum_j\pi_jK_j(a_j^*).
\end{align*}
For two regimes and a fixed tie rule put
$A=\sum_j\pi_j[K_j(a_{1-j}^*)-K_j(a_j^*)]$.
This nonnegative quantity bounds the leading cost of a wrong commitment.
The model laws, prior, initialization radius and probe schedule are public;
the realized regime and initial error are not.

\begin{theorem}[Paid separation from every fixed training prefix]
Suppose a deterministic, target-invariant Gaussian probe has length $n$,
each directional likelihood error is at most $0<\delta<1/2$, and the continuing
programs commit to $a_j^*$ on its decision.
If $\Delta>\delta A$, its signed same-stream upper certificate is below the
minimum worst-initialization prior risk of every full-budget fixed program
and every observation-independent randomized mixture for all sufficiently
large $B$.
More precisely the certified advantage is at least
\[
\frac{\Delta-\delta A}{B}
-\frac{C(n+1)}{B^2},
\]
for a finite model/catalogue-dependent $C$, once all post-probe horizons are
positive fractions of $B$.
If instead the regimes share a leading-coefficient minimizing action, then
$\Delta=0$, and no selector from the maximal-affordable-horizon post-probe
catalogue plus stop, following any such probe with $n=O(\log B)$, can
improve on the best fixed prefix by more than $O((n+1)/B^2)$.
\label{thm:research-leading-separation}
\end{theorem}

The positive claim is relative to a fixed learner and a complete resource-and-
prefix class, not to arbitrary estimators or learning-rate sequences.
The negative claim is an order statement for a specified sublinear paid probe;
it is not an impossibility theorem over all sequential sensing policies.

\subsection{Proof and the role of every hypothesis}
For a fixed action and $N$ applied observations, the exact running-mean risk is
\[
R_{j,a}(N)=\frac{r^2+N/q+
\theta_j[N+2\sum_{s=1}^{N-1}(N-s)\lambda^{bs}]}{(N+1)^2}.
\]
The geometric-sum formula gives
$R_{j,a}(N)=v_j(a)/N+O(N^{-2})$, uniformly on the finite catalogue.
It also gives
$R_{j,a}(N+1)-R_{j,a}(N)=-v_j(a)/N^2+O(N^{-3})$
with an exponentially small remainder.
Since $v_j(a)>0$, each action's risk, and each prior average of it, decreases
strictly beyond a fixed horizon.
Every finite early prefix has positive risk, whereas its maximal affordable
prefix has risk tending to zero.
Thus the optimum over all prefixes eventually uses each candidate action's
maximal horizon, and its global optimum is
$I_{\rm full}=K_\pi/B+O(B^{-2})$.
The same endpoint initialization maximizes all scalar fixed risks; randomizing
fixed programs cannot lower their minimum prior risk.

Subtracting a length-$n$ paid probe changes each maximal horizon from
$a_aB+O(1)$ to $a_aB+O(n+1)$, including rounding and bounded withheld updates.
The correct-program prior risk is therefore
$\sum_j\pi_jK_j(a_j^*)/B+O((n+1)/B^2)$.
The directional-error term is at most $\delta A/B+O((n+1)/B^2)$.
The even invariant likelihood decision cancels the unknown initialization's
odd term; the signed continuing-state correction is $O(B^{-2})$, because
$k_a=O(B^{-1})$ and the probe's explained variance is at most $\theta_j$.
The constants in its truncated-normal factors are finite at fixed $\delta$.
Combining these facts with the fixed-prefix expansion proves the lower bound
on the certified advantage.

For the negative claim, a common leading minimizer gives
$K_\pi=\sum_j\pi_j\min_aK_j(a)$.
The post-probe perfect-regime fixed risk consequently has that same leading
coefficient and a remainder $O((n+1)/B^2)$.
The preceding oracle-plus-half-state-ceiling bound applies to every measurable
randomized selector, whose state-value ceiling is $O(B^{-2})$.
Subtracting the fixed-prefix risk proves the stated upper order.
No monotonicity of an averaged affine learner is invoked in this running-mean
theorem.

\subsection{Existence of an observable finite-confidence probe}
For $\Sigma_j(n)=\theta_jR_\lambda(n)+I/Q$, with distinct positive variances,
the eigenvalues of $R_\lambda(n)$ lie in a fixed positive compact interval.
Both covariance matrices are functions of the same $R_\lambda(n)$.
Each eigen-coordinate contributes a Gaussian affinity factor strictly below
one, uniformly on that interval.
Thus the full centered-noise affinity is at most $\kappa^n$ for some
$\kappa<1$.
The contrast determinant bound gives $\rho_D(n)\le C_1\kappa^n$ for the
unknown scalar target.
Any fixed confidence $\delta>0$ therefore has a finite, observable paid probe.
Its actual directional bound is recomputed after projection, not inherited
from the centered-noise experiment.

\subsection{An explicit sufficient budget threshold}
The sufficiently-large-budget statement has computable, albeit conservative, constants.
All action maxima in this subsection range over non-stop resource pairs; the zero-update program is treated separately.
Write $\rho_a=\lambda^{b_a}$ and define
\[
c_{j,a}=r^2+2v_j(a)+\frac{2\theta_j\rho_a}{(1-\rho_a)^2}.
\]
The geometric sum gives the exact identity
\[
R_{j,a}(N)=\frac{r^2+Nv_j(a)-2\theta_j\rho_a(1-\rho_a^N)/(1-\rho_a)^2}{(N+1)^2},
\]
and hence $|R_{j,a}(N)-v_j(a)/N|\le c_{j,a}/N^2$ for every $N\ge1$.
Let $p_m,p_e$ be the probe's full paid message and clock charges, including any fixed startup cost, and put
\[
L_a(p)=D+1+\max\{(1+p_m)/(h+q_a),(1+p_e)/b_a\}.
\]
For the declared FIFO convention, rounding the two budgets and subtracting the $D$ withheld updates gives
$a_aB-L_a(p)\le N_a^{\rm post}\le a_aB$.
Thus $B\ge2L_a(p)/a_a$ implies
\begin{align*}
\left|J_{j,a}^{r}-\frac{K_j(a)}B\right|&\le\frac{C_{j,a}(p)}{B^2},\\
C_{j,a}(p)&=\frac{4c_{j,a}+2v_j(a)L_a(p)}{a_a^2}.
\end{align*}

The all-prefix lower bound can also avoid an unspecified eventual-monotonicity threshold.
For each action use the prior averages $\bar v_a=\sum_j\pi_jv_j(a)$ and $\bar c_a=\sum_j\pi_jc_{j,a}$, and set $n_{0,a}=\lceil2\bar c_a/\bar v_a\rceil$.
For $N\ge n_{0,a}$ the expansion lower bound is at least $\bar v_a/(2N)$.
If $N<N_{\max,a}/4$ this exceeds $\bar v_a/N_{\max,a}$; if $N\ge N_{\max,a}/4$, the expansion instead bounds the remainder by $16\bar c_a/N_{\max,a}^2$.
For $N<n_{0,a}$, including stop, the initialization term is at least $r^2/(n_{0,a}+1)^2$.
Consequently every prefix obeys
\[
\bar R_a(N)\ge\frac{\bar v_a}{N_{\max,a}}-\frac{16\bar c_a}{N_{\max,a}^2}
\]
whenever $N_{\max,a}\ge\bar v_a(n_{0,a}+1)^2/r^2$.
It follows that
\[
I_{\rm full}\ge\frac{K_\pi}B-\frac{C_F}{B^2},\quad
C_F=\max_a\frac{64\bar c_a}{a_a^2},
\]
provided $B\ge B_F$, where
\[
B_F=\max\left\{\frac{K_\pi}{r^2},\max_a\frac{2L_a(0)}{a_a},
\max_a\frac{2\bar v_a(n_{0,a}+1)^2}{a_ar^2}\right\}.
\]
The zero-update program has risk $r^2$ and is covered by the first term.

Let $c_j=a_j^*$, $w_j=a_{1-j}^*$, and $\kappa=\max_a4\rho_a^2/[a_a^2(1-\rho_a)^2]$.
At $B\ge\max_a2L_a(p)/a_a$, each common-state coefficient has $k_a^2\le\kappa/B^2$.
The explained variance is at most $\theta_j$, and the even decision has centered-square selection correction at most $4\varphi(1)\theta_j\kappa/B^2$.
This also bounds the corresponding signed truncated-normal correction from above.
Using the directional error bound $\delta$ and $K_j(w_j)\ge K_j(c_j)$ therefore gives a certified upper risk
\[
U\le\frac{\sum_j\pi_jK_j(c_j)+\delta A}{B}+\frac{C_U}{B^2},
\]
where
\begin{align*}
C_U=\sum_j\pi_j\{&C_{j,c_j}(p)+\delta[C_{j,c_j}(p)+C_{j,w_j}(p)]\\
&+4\varphi(1)\theta_j\kappa\}.
\end{align*}
Thus a completely explicit sufficient threshold is
\[
B\ge\max\{B_F,\max_a2L_a(p)/a_a\},\quad
B>\frac{C_F+C_U}{\Delta-\delta A},
\]
when $\Delta>\delta A$.
All constants use public model, catalogue and paid-probe quantities.
They may be much more conservative than evaluating the exact finite-budget certificate; no small practical threshold or observed improvement is inferred from this construction.
For a length-$n$ probe with bounded per-observation charges, $C_U=O(n+1)$, recovering the order in the theorem.

\subsection{A rational resource witness on the existing public grid}
Take $Q=8$, $h=8$, $\theta_0=1/20$, $\theta_1=2$, $\lambda=1/5$,
balanced prior, $\beta_m=16$, $\beta_e=4$, and the original
$q\in\{1,2,4,8\}$, $b\in\{1,2,4,8\}$ catalogue.
Exact enumeration of its sixteen leading coefficients yields
\begin{align*}
a_0^*&=(8,4),\quad K_0(a_0^*)=1093/6240,\\
a_1^*&=(2,2),\quad K_1(a_1^*)=5/3,
\end{align*}
while the best balanced fixed coefficient is $K_\pi=773/768$ at $(2,2)$.
The wrong-commitment coefficient is $A=407/1280$.
At directional bound $\delta=1/40$, the sufficient leading margin
$\Delta-\delta A$ is strictly positive (approximately $0.07765$).
Its ratio to $K_\pi$ is approximately $7.71$ percent.
This is a model-population, sufficiently-large-budget witness, not a finite-
grid empirical gain or a claimed threshold value.
All finite budget thresholds and admitted finite-grid designs must still be
computed from the actual paid projected probe and exact learner risk.
